\documentclass[11pt]{article}
\usepackage[a4paper,margin=27mm]{geometry}
\usepackage{amsmath,amssymb,mathtools}
\usepackage{fontspec,unicode-math}
\setmainfont{Latin Modern Roman}
\setmathfont{Latin Modern Math}
\usepackage{xurl}
\usepackage[hidelinks]{hyperref}
\newcommand{\sref}[2]{\hyperlink{source:#1:#2}{[S#2]}}
\newcommand{\cataloguescope}[1]{\paragraph{Mathematical question.}#1}
\hypersetup{pdftitle={UnsolvedMath Problem 9400168}}
\begin{document}
% UnsolvedMath ID 9400168; source catalogue code AMR-093-0168
\setcounter{section}{9400167}
\section{\texorpdfstring{Can a prime p satisfy $2^{p-1}\equiv 1\pmod{p^2}$ and $3^{p-1}\equiv 1\pmod{p^2}$ simultaneously}{Can a prime p satisfy 2\textasciicircum{}\{p-1\}congruent to 1pmod\{p\textasciicircum{}2\} and 3\textasciicircum{}\{p-1\}congruent to 1pmod\{p\textasciicircum{}2\} simultaneously}}\label{9400168}
{\small\sffamily UnsolvedMath ID 9400168 · Number Theory}
\subsection{Definitions and mathematical statement}

\paragraph{Shared notation.}$\mathbb N=\{1,2,\ldots\}$, and $\mathbb Z,\mathbb Q,\mathbb R,\mathbb C$ denote the integers, rationals, reals and complex numbers. Any other conventions are specified locally.


For a prime $p\nmid a$, being Wieferich to base $a$ means $a^{p-1}\equiv1\pmod{p^2}$. A prime satisfying this for both bases $2$ and $3$ must be different from $2$ and $3$.
\paragraph{Statement recorded in the supplied source.}
Does there exist a prime $p$ satisfying both $2^{p-1}\equiv1\pmod{p^2}$ and $3^{p-1}\equiv1\pmod{p^2}$ simultaneously? \sref{9400168}{2}
\subsection{Short English statement}
Can one prime be Wieferich for both bases two and three?
\subsection{Sources}
\begin{description}
\item[\hypertarget{source:9400168:1}{S1}] ULAM.AI and the UnsolvedMath Contributors, UnsolvedMath: A Curated Collection of Open Mathematics Problems, record 9400168 (AMR-093-0168). CC BY 4.0. \url{https://huggingface.co/datasets/ulamai/UnsolvedMath}.
\item[\hypertarget{source:9400168:2}{S2}] Wikipedia contributors, \emph{List of unsolved problems in mathematics}, revision 1366636767 (2026) Number theory, Can a prime p satisfy $2^{p-1}\equiv 1\pmod{p^2}$ and $3^{p-1}\equiv 1\pmod{p^2}$ simultaneously. \url{https://en.wikipedia.org/w/index.php?title=List_of_unsolved_problems_in_mathematics&oldid=1366636767#Number_theory}.
\item[\hypertarget{source:9400168:3}{S3}] Jonathan Sondow, \emph{Lerch quotients, Lerch primes, Fermat--Wilson quotients, and the Wieferich-non-Wilson primes $2,3,14771$}, arXiv:1110.3113v5 (2012) Introduction and Section 3.2, pp. 1--2 and 9--10. \url{https://arxiv.org/pdf/1110.3113}.
\end{description}
\label{end:9400168}
\end{document}
